\documentclass[11pt]{article}

\usepackage[margin=1.02in]{geometry}
\usepackage{amsmath,amssymb,amsthm,amscd,mathtools}
\usepackage{booktabs}
\usepackage{array}
\usepackage{enumitem}
\usepackage{microtype}
\usepackage{xcolor}
\usepackage[hidelinks]{hyperref}
\usepackage{fancyhdr}

\definecolor{deepblue}{HTML}{173B57}
\hypersetup{
  colorlinks=true,
  linkcolor=deepblue,
  citecolor=deepblue,
  urlcolor=deepblue,
  pdftitle={An exceptional four-dimensional unitary Hecke Yang--Baxter operator: a five-word Pauli--Clifford normal form},
  pdfauthor={Alec Kriebel}
}

\pagestyle{fancy}
\fancyhf{}
\fancyhead[L]{\small\textsc{Exceptional Hecke localization}}
\fancyhead[R]{\small\textsc{A. Kriebel}}
\fancyfoot[C]{\thepage}
\setlength{\headheight}{14pt}

\newtheorem{theorem}{Theorem}[section]
\newtheorem{lemma}[theorem]{Lemma}
\newtheorem{proposition}[theorem]{Proposition}
\newtheorem{corollary}[theorem]{Corollary}
\theoremstyle{remark}
\newtheorem{remark}[theorem]{Remark}

\newcommand{\Id}{\mathbf 1}
\newcommand{\Tr}{\operatorname{Tr}}
\newcommand{\End}{\operatorname{End}}
\newcommand{\rank}{\operatorname{rank}}
\newcommand{\bbC}{\mathbb C}
\newcommand{\bbQ}{\mathbb Q}
\newcommand{\cC}{\mathcal C}
\newcommand{\cR}{\mathcal R}
\newcommand{\boxt}{\mathbin{\boxtimes}}

\title{\vspace{-1.4cm}\textbf{An exceptional four-dimensional unitary}\\[0.2em]
\textbf{Hecke Yang--Baxter operator:}\\[0.2em]
\textbf{a five-word Pauli--Clifford normal form}}
\author{Alec Kriebel\\[-0.1em]
\small Independent researcher\\[-0.1em]
\small \href{mailto:me@aleckriebel.com}{me@aleckriebel.com}
\quad\textperiodcentered\quad
\href{https://orcid.org/0009-0001-9320-500X}{ORCID 0009-0001-9320-500X}}
\date{Version 1.2.0 --- 19 August 2026}

\begin{document}
\maketitle

\begin{abstract}
We give an explicit \(16\times16\) unitary Yang--Baxter operator in the class
\[
  \left[e^{i\pi/3},\frac12,4\right].
\]
On \(\bbC^4\cong\bbC^2\otimes\bbC^2\), its reflection has a five-word real
Pauli--Clifford expansion.  Four words combine into an involution \(M\), a fifth
involution \(E\) anticommutes with \(M\), and an explicit \(18\)-word cubic
certificate restricts the reflection circle \(\alpha M+\beta E\) exactly to
\[
  \alpha^2=\frac23,\qquad \beta^2=\frac13.
\]
The spectral projection has both unnormalized partial traces \(2\Id_4\), so
normalized matrix trace induces the \(\eta=1/2\) Markov trace and a faithful
representation of \(H_n(3,6)\) for every \(n\).  A direct dimension-three
obstruction, together with the known dimension-two obstruction, proves that
four is minimal.  Galindo and Rowell independently obtained the same local
solution through a quaternionic twisted-group-algebra construction.  We
identify the present matrix exactly with a local-unitary normal form of the
opposite of their Family~III operator and extend the comparison to all braid
strands.  Transporting known Family~III results shows that the braid images are
finite and Clifford in a fixed conjugated Pauli frame.  A direct scalar
enhancement and skein calculation identify the trace invariant as
\(2P_{\mathrm H}(L;i,i)\); the Galindo--Rowell metric-family algorithm then
gives exact deterministic polynomial-time evaluation.
\end{abstract}

\medskip
\noindent\textbf{Keywords.}
Yang--Baxter operator; Hecke algebra; braid group representation; unitary
localization; fusion category; link invariant; Clifford group.

\smallskip
\noindent\textbf{2020 Mathematics Subject Classification.}
20C08 (primary); 20F36, 16T25, 57K14, 81R50 (secondary).

\section{Introduction and main result}

Let \(V\) be a finite-dimensional complex Hilbert space.  A unitary
\(R\in\End(V\otimes V)\) is a Yang--Baxter operator when
\begin{equation}
  (R\otimes\Id_V)(\Id_V\otimes R)(R\otimes\Id_V)
  =
  (\Id_V\otimes R)(R\otimes\Id_V)(\Id_V\otimes R).
  \label{eq:ybe}
\end{equation}
Lechner~\cite{Lechner2026} classified the possible characters of
non-involutive unitary \(R\)-matrices with two eigenvalues, up to one
exceptional family.  In his normalization the two eigenvalues are
\(-1\) and \(q\), the \((-1)\)-spectral projection \(P\) has normalized
trace
\[
  \eta=\frac{\rank P}{(\dim V)^2},
\]
and the classes left unresolved by that analysis are
\[
  \left[e^{i\pi/3},\frac12,2m\right]\qquad(m\ge1).
\]
The class in base dimension two is empty; base dimension four was therefore
the smallest unresolved member of this remaining family.

The ordinary-localization framework and conjecture were introduced by Rowell
and Wang~\cite{RowellWang2012}.  The specific question here was highlighted
in the case study of Galindo, Hong, and Rowell~\cite[\S5.6]{GHR2013}.  They
showed that the Jones--Wenzl
representations associated with \(\cC(\mathfrak{sl}_3,6)\) have a
two-dimensional unitary quasi-localization and an \(8\times8\)
\((3,1)\)-generalized localization,
but no two-dimensional ordinary unitary localization.  They wrote that the
sequence did not appear to have an ordinary localization and noted it as a
possible counterexample to the Rowell--Wang localization conjecture, restated
as GHR Conjecture~1.5.  Lechner isolated
this family as the sole remaining case in his classification preprint and
explicitly identified it with the unknown localization arising from the GHR
\(\cC(\mathfrak{sl}_3,6)\) case study
\cite[p.~16, discussion following Proposition~3.6]{Lechner2026}.  The present
paper supplies an explicit operator in that family.  The dimension-four
construction and localization proof are independent of Lechner's
classification; the minimality argument in Section~\ref{sec:consequences}
uses \cite[Lemma~3.1 and Theorem~3.4]{Lechner2026}.

\medskip
\noindent\textbf{Note added: independent concurrent work.}
Version 1.1.0 of this manuscript, containing the explicit operator and the
existence, localization, and minimum-dimension results, was publicly released
on 28 July 2026.\footnote{Release tag
\href{https://github.com/AlecKriebel/Math/releases/tag/exceptional-ybe-d4-v1.1.0}
{\texttt{exceptional-ybe-d4-v1.1.0}}, published at 2026-07-28 04:10:58 UTC;
artifact commit
\texttt{e2669c5b2f99\allowbreak338c79381dc4\allowbreak2bdbc61ee8b9\allowbreak63c3};
historical
PDF SHA-256
\texttt{af4ff57c4b8c5cd3\allowbreak7f47f8a6da880b4f\allowbreak93b9c22d6e2908a3\allowbreak{}ef1f6ebf5fb1d049}.}
Galindo and Rowell subsequently posted an independent quaternionic
construction~\cite{GalindoRowell2026} and report having obtained and privately
circulated it earlier.  The works are treated here as independent and
concurrent; no claim is made about private discovery priority.
Section~\ref{sec:quaternionic} compares the local formulas exactly, and
Section~\ref{sec:global} transfers the resulting braid and link structures to
the five-word representative.

The following is the central result of this paper.

\begin{theorem}\label{thm:main}
Let
\[
  q=e^{i\pi/3}=\frac{1+i\sqrt3}{2}.
\]
There is an explicit unitary \(R\in\End(\bbC^4\otimes\bbC^4)\) such that
\[
  (R+\Id_{16})(R-q\Id_{16})=0,\qquad
  R_1R_2R_1=R_2R_1R_2,
\]
and each eigenvalue \(-1,q\) has multiplicity eight.  With the conventions
\(\Tr_1=\Tr\otimes\operatorname{id}\) and
\(\Tr_2=\operatorname{id}\otimes\Tr\), its
\((-1)\)-spectral projection \(P\) satisfies
\[
  \Tr_1(P)=\Tr_2(P)=2\Id_4.
\]
Therefore
\[
  \left[e^{i\pi/3},\frac12,4\right]\ne\varnothing,
\]
and the associated tensor-space representations induce faithful
embeddings
\[
  H_n(3,6)\hookrightarrow\End\!\left((\bbC^4)^{\otimes n}\right)
  \qquad(n\ge2).
\]
Under the identification in~\cite[\S5.6]{GHR2013}, these embeddings form an
ordinary unitary localization of the
\(\cC(\mathfrak{sl}_3,6)\) Jones--Wenzl representation sequence.
\end{theorem}

\begin{corollary}\label{cor:minimal}
The minimum local dimension of an ordinary unitary localization of this
Jones--Wenzl sequence is four.
\end{corollary}

\begin{theorem}[Global braid and link consequences]\label{thm:global-summary}
Put \(\kappa=q-1=q^2=e^{2\pi i/3}\), and let
\(\rho_n:B_n\to\mathrm U((\bbC^4)^{\otimes n})\) be generated by the adjacent copies of
\(R\).  Then \((R,\Id_4,\kappa,2)\) is an enhanced Yang--Baxter operator, with
oriented-link invariant
\[
  \mathcal J_R(\widehat\xi)
  =\kappa^{-\operatorname{wr}(\xi)}2^{-n}\Tr(\rho_n(\xi)).
\]
For the normalized HOMFLYPT polynomial used below,
\[
  \mathcal J_R(L)=2P_{\mathrm H}(L;i,i).
\]
If \(c(L)\) is the number of components and
\(d_2(L)=\dim_{\mathbb F_2}H_1(\Sigma_3(L);\mathbb F_2)\) for the oriented
triple cyclic cover specified in Section~\ref{sec:global}, then
\[
  \mathcal J_R(L)=2(-1)^{c(L)-1}(-2)^{d_2(L)/2}.
\]
Via the all-strand equivalence of Theorem~\ref{thm:all-strand} and the
Family~III results cited in Section~\ref{sec:global}, every image
\(\rho_n(B_n)\) is finite, its generators are Clifford operators in a fixed
conjugated Pauli frame, and \(\mathcal J_R\) is exactly computable by a
deterministic classical algorithm in time polynomial in the strand number and
braid-word length.
\end{theorem}

The formula for \(R\) is given in Section~\ref{sec:construction}; the matrix
identities are proved in Sections~\ref{sec:reflection} and
\ref{sec:ybe}.  Section~\ref{sec:localization} proves the full localization
claim rather than only the two-site spectral statement.  The minimality
argument is in Section~\ref{sec:consequences}.  Section~\ref{sec:gybe}
records a generalized Yang--Baxter normal form, Section~\ref{sec:quaternionic}
gives the quaternionic factorization and concurrent comparison, and
Section~\ref{sec:global} proves Theorem~\ref{thm:global-summary}.
Section~\ref{sec:verification} records verification and provenance, and
Section~\ref{sec:priority} discusses related work, novelty, and limitations.

\section{The five-word construction}\label{sec:construction}

Put \(V=\bbC^2\otimes\bbC^2\cong\bbC^4\), and define the real matrices
\[
  I=\begin{pmatrix}1&0\\0&1\end{pmatrix},\quad
  X=\begin{pmatrix}0&1\\1&0\end{pmatrix},\quad
  Z=\begin{pmatrix}1&0\\0&-1\end{pmatrix},\quad
  J=\begin{pmatrix}0&-1\\1&0\end{pmatrix}.
\]
Here \(J=XZ=-iY\), so \(\{I,X,Z,J\}\) is the real Pauli--Clifford basis of
\(M_2(\bbC)\); every tensor word in this basis is a real signed orthogonal
matrix.
A word such as \(ZIZZ\) denotes
\(Z\otimes I\otimes Z\otimes Z\) on \(V\otimes V\).  Define
\begin{equation}
\begin{split}
  H={}&-\frac1{\sqrt6}ZIZZ-\frac1{\sqrt6}ZIJJ
        -\frac1{\sqrt6}JIZJ+\frac1{\sqrt6}JIJZ\\
      &-\frac1{\sqrt3}XIXX,
  \label{eq:H}
\end{split}
\end{equation}
and
\begin{equation}
  R=\frac{q-1}{2}\Id_{16}+\frac{q+1}{2}H
   =\frac{-1+i\sqrt3}{4}\Id_{16}
    +\frac{3+i\sqrt3}{4}H.
  \label{eq:R}
\end{equation}
Equations~\eqref{eq:H}--\eqref{eq:R} are an exact specification of all
\(256\) matrix entries in the computational basis.
The visible identity in the second qubit is not an ordinary two-site identity
amplification of a \(4\times4\) Yang--Baxter operator: the active part uses one
qubit of the first four-dimensional site and both qubits of the second.
Section~\ref{sec:gybe} identifies it instead as a \((3,2)\)-generalized
operator.

The relevant multiplication table is
\[
  X^2=Z^2=I,\qquad J^2=-I,
\]
\[
  XZ=J=-ZX,\qquad XJ=Z=-JX,\qquad JZ=X=-ZJ.
\]
Thus every product of tensor words reduces to a signed word over
\(\{I,X,Z,J\}\).

\section{A circle of reflections}\label{sec:reflection}

The five terms in~\eqref{eq:H} have a compact algebraic explanation.  Set
\[
  A=ZIZZ,\quad B=ZIJJ,\quad C=JIZJ,\quad D=JIJZ,\quad E=XIXX
\]
and
\begin{equation}
  M=\frac{-A-B-C+D}{2}.
  \label{eq:M}
\end{equation}

\begin{lemma}\label{lem:reflection}
The operators \(A,B,C,D,E\) are real symmetric involutions.
The first four commute, \(D=ABC\), and \(E\) anticommutes with each of
\(A,B,C,D\).  Consequently
\[
  M^2=E^2=\Id_{16},\qquad ME=-EM.
\]
For real \(\alpha,\beta\),
\[
  H(\alpha,\beta)=\alpha M+\beta E
\]
satisfies
\[
  H(\alpha,\beta)^2=(\alpha^2+\beta^2)\Id_{16}.
\]
\end{lemma}

\begin{proof}
The word multiplication table gives every stated commutation relation and
\(D=ABC\).  One may also prove \(M^2=\Id\) without expanding.  Simultaneously
diagonalize \(A,B,C\), and let their joint eigenvalues be
\(x,y,z\in\{\pm1\}\).  Since \(D=ABC\), the eigenvalue of \(M\) is
\[
  \frac{-x-y-z+xyz}{2}\in\{\pm1\}.
\]
The anticommutation \(ME=-EM\) then makes the cross term in
\((\alpha M+\beta E)^2\) vanish.
\end{proof}

\begin{remark}[Graph-phase form]
If \(x=(-1)^a,y=(-1)^b,z=(-1)^c\), then
\[
  \frac{-x-y-z+xyz}{2}=-(-1)^{ab+ac+bc}.
\]
Thus in the joint stabilizer eigenbasis, \(M\) is the negative phase
associated with the triangle graph, while \(E\) flips all three syndrome
bits.  This explains the four-plus-one support in~\eqref{eq:H}; no graph-state
formalism is needed in the proof.
\end{remark}

\section{The cubic certificate and Yang--Baxter equation}\label{sec:ybe}

For any \(K\in\End(V\otimes V)\), write
\[
  K_1=K\otimes\Id_V,\qquad K_2=\Id_V\otimes K.
\]
The Yang--Baxter equation for~\eqref{eq:R} will follow from a cubic identity
for \(H\).

\begin{proposition}\label{prop:circle}
On the unit circle \(\alpha^2+\beta^2=1\),
\begin{equation}
  H(\alpha,\beta)_1H(\alpha,\beta)_2H(\alpha,\beta)_1
  -
  H(\alpha,\beta)_2H(\alpha,\beta)_1H(\alpha,\beta)_2
  =
  \frac13\bigl(H(\alpha,\beta)_1-H(\alpha,\beta)_2\bigr)
  \label{eq:cubic-family}
\end{equation}
if and only if
\[
  \beta^2=\frac13,\qquad\alpha^2=\frac23.
\]
In particular, the four signed points with these squared coordinates are
exactly the solutions on this reflection circle.
\end{proposition}

\begin{proof}
Reduce the left side of~\eqref{eq:cubic-family}, minus its right side, in
the six-letter word basis.  Exactly the following \(18\) words survive:
\begin{center}
\small
\renewcommand{\arraystretch}{1.16}
\begin{tabular}{>{\raggedright\arraybackslash}p{0.33\textwidth}
                >{\raggedright\arraybackslash}p{0.59\textwidth}}
\toprule
coefficient & words\\
\midrule
\(-\alpha(3\beta^2-1)/6\) &
\(IIJIJZ,\ JIZJII,\ ZIJJII,\ ZIZZII\)\\
\(+\alpha(3\beta^2-1)/6\) &
\(IIJIZJ,\ IIZIJJ,\ IIZIZZ,\ JIJZII\)\\
\(-\beta(3\alpha^2-3\beta^2-1)/3\) &
\(IIXIXX\)\\
\(+\beta(3\alpha^2-3\beta^2-1)/3\) &
\(XIXXII\)\\
\(+\alpha(\alpha^2-2\beta^2)/2\) &
\(JIJZXX,\ JIZJXX,\ XIJXZJ,\ ZIZZXX\)\\
\(-\alpha(\alpha^2-2\beta^2)/2\) &
\(XIJXJZ,\ XIZXJJ,\ XIZXZZ,\ ZIJJXX\)\\
\bottomrule
\end{tabular}
\end{center}
This is a finite certificate: each row follows by repeated use of the
displayed \(2\times2\) multiplication table.  Since \(\{I,X,Z,J\}\) is a
basis of \(M_2(\bbC)\), its sixfold tensor words form a basis of
\(\End((\bbC^2)^{\otimes6})\).  Hence the residual vanishes if and only if
every displayed coefficient vanishes.  On
\(\alpha^2+\beta^2=1\),
\[
  \alpha^2-2\beta^2=1-3\beta^2,\qquad
  3\alpha^2-3\beta^2-1=2(1-3\beta^2).
\]
Hence all coefficients vanish when \(3\beta^2=1\).  Conversely, if
\(\alpha=0\), the coefficient of \(IIXIXX\) is
\(4\beta/3\ne0\).  If \(\beta=0\), the coefficient of \(IIJIJZ\) is
\(\alpha/6\ne0\).  If \(\alpha\beta\ne0\), the first coefficient group forces
\(3\beta^2=1\), and the unit-circle condition then gives
\(\alpha^2=2/3\).  The claim follows.
\end{proof}

Choosing
\[
  \alpha=\sqrt{\frac23},\qquad\beta=-\frac1{\sqrt3}
\]
in \(H(\alpha,\beta)\) gives exactly~\eqref{eq:H}.  Lemma~\ref{lem:reflection}
and Proposition~\ref{prop:circle} therefore give
\begin{equation}
  H^*=H,\qquad H^2=\Id_{16},\qquad
  H_1H_2H_1-H_2H_1H_2=\frac13(H_1-H_2).
  \label{eq:H-identities}
\end{equation}
Every word in~\eqref{eq:H} contains a traceless factor, so
\(\Tr H=0\).  The eigenvalues \(+1\) and \(-1\) thus both have multiplicity
eight.

\begin{proof}[Proof of the matrix part of Theorem~\ref{thm:main}]
Put
\[
  a=\frac{q-1}{2},\qquad b=\frac{q+1}{2},
\]
so \(R=a\Id+bH\).  Since \(H_i^2=\Id\), direct expansion gives
\[
  R_1R_2R_1-R_2R_1R_2
  =
  b\left[
    a^2(H_1-H_2)+b^2(H_1H_2H_1-H_2H_1H_2)
  \right].
\]
But
\[
  \frac ab=\frac{q-1}{q+1}=\frac{i}{\sqrt3},
\]
so \(a^2+b^2/3=0\).  Equation~\eqref{eq:H-identities} proves
the Yang--Baxter equation.

On the \(+1\)-eigenspace of \(H\), \(R=a+b=q\); on the
\(-1\)-eigenspace, \(R=a-b=-1\).  Therefore
\[
  (R+\Id)(R-q\Id)=0,\qquad R^*R=\Id,
\]
with both eigenvalues of multiplicity eight.  The \((-1)\)-spectral
projection is
\[
  P=\frac{\Id_{16}-H}{2},
\]
so \(P=P^*=P^2\), \(\rank P=8\), and \(\eta=8/16=1/2\).
Finally, each word in \(H\) has a traceless factor on each four-dimensional
site.  Hence
\[
  \Tr_1(H)=\Tr_2(H)=0,\qquad
  \Tr_1(P)=\Tr_2(P)=2\Id_4.
\]
\end{proof}

\section{The full localization statement}\label{sec:localization}

Galindo--Rowell also establish strict localization in local dimension four
\cite[\S13]{GalindoRowell2026}.
This section and Section~\ref{sec:consequences} retain the independent direct
proofs in public version 1.1.0, including scalar partial traces, the
kernel--annihilator equality, and the explicit dimension-three obstruction.

We now verify that the construction localizes the whole Jones--Wenzl tower,
not merely its two-site generator.  Let \(H_n(q)\) be the specialized Hecke
algebra of type \(A_{n-1}\), generated by \(g_1,\ldots,g_{n-1}\), with the braid and
far-commutativity relations and
\[
  (g_i+\Id)(g_i-q\Id)=0.
\]
Its idempotent generators \(e_i\) are related to the braid generators by
\[
  g_i=q\Id-(1+q)e_i.
\]
The Yang--Baxter and Hecke relations define a representation
\[
  \rho_n:H_n(q)\longrightarrow\End(V^{\otimes n}),\qquad
  \rho_n(g_i)=R_i,\qquad \rho_n(e_i)=P_i.
\]
Let
\[
  \tau_n(T)=4^{-n}\Tr(T)
\]
be normalized matrix trace.

Equip \(H_n(q)\) with the conjugate-linear involution
\[
  e_i^*=e_i,\qquad\text{equivalently}\qquad g_i^*=g_i^{-1}.
\]
Since \(P_i=P_i^*\) and \(R_i\) is unitary, each \(\rho_n\) is a unital
\(*\)-representation.  If
\(\iota_n:H_n(q)\to H_{n+1}(q)\) is the standard inclusion, then
\begin{equation}
  \rho_{n+1}(\iota_n(x))=\rho_n(x)\otimes\Id_4.
  \label{eq:tower-compatibility}
\end{equation}
For the \(\eta=1/2\) Markov trace, write
\[
  \operatorname{Ann}_n
  =
  \left\{x\in H_n(q):
  \operatorname{tr}^{(n)}_{1/2}(yx)=0\ \text{for every }y\in H_n(q)\right\}.
\]

\begin{proposition}\label{prop:markov}
The traces \(\tau_n\circ\rho_n\) are the Markov traces with parameter
\(\eta=1/2\):
\begin{equation}
  \tau_{n+1}\!\left((\rho_n(x)\otimes\Id_4)P_n\right)
  =\frac12\,\tau_n(\rho_n(x))
  \qquad(x\in H_n(q)).
  \label{eq:markov}
\end{equation}
For every \(n\),
\[
  \ker\rho_n=\operatorname{Ann}_n.
\]
Consequently \(\rho_n\)
induces an injective representation of the semisimple quotient
\[
  H_n(3,6)=H_n(q)/\operatorname{Ann}_n.
\]
\end{proposition}

\begin{proof}
Taking the partial trace over the last tensor factor and using
\(\Tr_2(P)=2\Id_4\) gives
\[
\begin{aligned}
  \tau_{n+1}\!\left((\rho_n(x)\otimes\Id_4)P_n\right)
  &=
  4^{-(n+1)}
  \Tr\!\left(\rho_n(x)\,\Tr_{n+1}(P_n)\right)\\
  &=
  4^{-(n+1)}\,2\Tr(\rho_n(x))
   =\frac12\tau_n(\rho_n(x)).
\end{aligned}
\]
Moreover,
\[
  \tau_{n+1}(\rho_n(x)\otimes\Id_4)=\tau_n(\rho_n(x)),
\]
so the traces are consistent under the standard inclusions.
Together with normalization and traciality, this is the defining Markov
property.  Uniqueness of the Markov trace with prescribed
\(\operatorname{tr}^{(n)}_{1/2}(e_i)=1/2\) identifies it with the trace defining
\(H_n(3,6)\)~\cite[Eq.~(3.2)]{Wenzl1988}.  Explicitly, in the normalization
used for this quotient,
\[
  \eta_{6,3}
  =\frac{1-q^{-2}}{(1+q)(1-q^{-3})}
  =\frac12,
  \qquad q=e^{i\pi/3}.
\]
This is the Wenzl Markov parameter for \(H_n(3,6)\).  GHR \(\S5.6\) identifies
this quotient tower with the \(\cC(\mathfrak{sl}_3,6)\) Jones--Wenzl sequence,
and Theorem~5.28 uses that identification for its generalized localization.

The equality \(\tau_n\circ\rho_n=\operatorname{tr}^{(n)}_{1/2}\) now gives
both kernel inclusions explicitly.  If \(x\in\operatorname{Ann}_n\), then
\[
\begin{aligned}
  x\in\operatorname{Ann}_n
  &\Longrightarrow
  0=\operatorname{tr}^{(n)}_{1/2}(x^*x)
   =\tau_n\!\left(\rho_n(x)^*\rho_n(x)\right)\\
  &\Longrightarrow \rho_n(x)=0,
\end{aligned}
\]
because normalized matrix trace is faithful.  Conversely, for every
\(y\in H_n(q)\),
\[
  \rho_n(x)=0
  \Longrightarrow
  \operatorname{tr}^{(n)}_{1/2}(yx)
  =\tau_n\!\left(\rho_n(y)\rho_n(x)\right)=0,
\]
so \(x\in\operatorname{Ann}_n\).  Thus the kernel is exactly the annihilator
and the induced map on the quotient is injective.
\end{proof}

\begin{remark}[Conventions in GHR \(\S5.6\)]
The published text of GHR \(\S5.6\) uses two \(q\)-normalizations: its opening
paragraph prints \(q=e^{2\pi i/3}\), whereas the braid eigenvalues and
Theorem~5.28 use \(q=e^{2\pi i/6}\).  We use the latter eigenvalue
normalization.  The proof of Theorem~5.28 also prints
\(\eta=(1-q^{-2})/(1+q^3)\), whose denominator vanishes at this value.  The
Wenzl parameter used in Proposition~\ref{prop:markov} is
\(\eta_{6,3}=1/2\).
\end{remark}

Equation~\eqref{eq:tower-compatibility} and Proposition~\ref{prop:markov}
give the quotient-tower compatibility explicitly:
\[
\begin{aligned}
  x\in\iota_n^{-1}(\operatorname{Ann}_{n+1})
  &\Longleftrightarrow \rho_{n+1}(\iota_n(x))=0\\
  &\Longleftrightarrow \rho_n(x)\otimes\Id_4=0\\
  &\Longleftrightarrow \rho_n(x)=0\\
  &\Longleftrightarrow x\in\operatorname{Ann}_n.
\end{aligned}
\]
Thus the inclusions descend to the semisimple quotients and give the
commuting square
\[
\begin{CD}
  H_n(3,6) @>{\overline\iota_n}>> H_{n+1}(3,6)\\
  @V{\overline\rho_n}VV @VV{\overline\rho_{n+1}}V\\
  \End(V^{\otimes n}) @>{T\mapsto T\otimes\Id_4}>>
  \End(V^{\otimes(n+1)}).
\end{CD}
\]

Galindo, Hong, and Rowell recall the braid-equivariant isomorphism
\[
  \bbC\,\rho_X(B_n)
  \cong\End_{\cC}(X^{\otimes n})\cong H_n(3,6)
\]
for the simple tensor generator \(X\) of \(\cC(\mathfrak{sl}_3,6)\), with
\(\operatorname{FPdim}(X)=2\)~\cite[\S5.6]{GHR2013}.  Since the braid image is
the entire algebra being localized, the compatible quotient embeddings above
are precisely the injective algebra maps required by
\cite[Definition~1.3]{GHR2013}.  Equivalently, by
\cite[Proposition~4.12 and Definition~4.13]{GHR2013}, they define an ordinary
unitary \(\mathrm{Hilb}_f\)-localization of the Jones--Wenzl sequence.

\begin{corollary}\label{cor:not-counterexample}
The Rowell--Wang localization conjecture
\cite[Conjecture~3.1, p.~601]{RowellWang2012}, as restated in
\cite[Conjecture~1.5]{GHR2013}, holds for the simple tensor generator \(X\) of
\(\cC(\mathfrak{sl}_3,6)\): one has
\[
  \operatorname{FPdim}(X)^2=4,
\]
and the associated Jones--Wenzl sequence admits an ordinary unitary
localization.  Hence the case singled out in
\cite[Remark~6.2(2)]{GHR2013} is not a counterexample.
\end{corollary}

\begin{proof}
Theorem~\ref{thm:main} supplies the asserted localization, and GHR \(\S5.6\)
gives \(\operatorname{FPdim}(X)=2\).  This does not contradict
\cite[Theorem~5.27]{GHR2013}: that result excludes a fiber functor by showing
that one would force a localization of dimension exactly
\(\operatorname{FPdim}(X)=2\), whereas the present ordinary localization has
local dimension four.
\end{proof}

\section{Minimality and amplification}\label{sec:consequences}

\begin{proof}[Proof of Corollary~\ref{cor:minimal}]
A one-dimensional local space cannot carry the required two-eigenvalue
generator.  Galindo, Hong, and Rowell prove that no two-dimensional ordinary
unitary localization exists~\cite[Lemma 5.26]{GHR2013}.  That lemma quotes the
two-dimensional classification up to an overall scalar: a representative
with spectrum \(\{-\chi,\chi q\}\) is normalized by multiplication by
\(\chi^{-1}\), which leaves its spectral projection and Markov parameter
unchanged.

It remains to exclude local dimension three.  Suppose such an ordinary
unitary localization existed, with local braid operator
\(F\in\End(\bbC^3\otimes\bbC^3)\).  Since a localization intertwines the
braid generators,
\[
  (F+\Id_9)(F-q\Id_9)=0.
\]
Injectivity at two strands rules out either spectral projection vanishing,
so both eigenvalues \(-1\) and \(q\) occur.  Thus \(F\) belongs to one of the
dimension-three Hecke classes in
\cite[Theorem~3.4]{Lechner2026}.  At \(q=e^{i\pi/3}\), the possible Markov
parameters are \(1/3,1/2,2/3\).  The value \(1/2\) is impossible in dimension
three because it would require a projection of rank \(9/2\).  By
\cite[Lemma~3.1]{Lechner2026}, the normalized matrix character of each such
Hecke operator is the corresponding positive Markov trace.

We now exclude the two remaining parameters directly.  In \(H_3(q)\), put
\[
  c=\frac{q}{(1+q)^2}=\frac13,
  \qquad T=e_1e_2e_1-ce_1.
\]
The projection-form Hecke relation is
\[
  e_1e_2e_1-e_2e_1e_2=c(e_1-e_2).
\]
For the Markov trace \(\operatorname{tr}^{(3)}_\eta\) with
\(\operatorname{tr}^{(3)}_\eta(e_i)=\eta\), multiply the displayed
projection-form Hecke relation on the right by \(e_2\).  Cyclicity and the
Markov property give
\[
\begin{aligned}
  \operatorname{tr}^{(3)}_\eta(e_1e_2e_1e_2)
  &=\operatorname{tr}^{(3)}_\eta(e_2e_1e_2)
    +c\operatorname{tr}^{(3)}_\eta(e_1e_2)
    -c\operatorname{tr}^{(3)}_\eta(e_2)\\
  &=\eta^2-c\eta+c\eta^2,
\end{aligned}
\]
while cyclicity and \(e_1^2=e_1\) give
\[
  \operatorname{tr}^{(3)}_\eta(e_1e_2e_1e_2e_1)
  =\operatorname{tr}^{(3)}_\eta(e_1e_2e_1e_2).
\]
Consequently,
\[
\begin{aligned}
  \operatorname{tr}^{(3)}_\eta(T^*T)
  &=\operatorname{tr}^{(3)}_\eta(e_1e_2e_1e_2)
    -2c\eta^2+c^2\eta\\
  &=(1-c)\eta(\eta-c).
\end{aligned}
\]
Thus a parameter-\(1/3\) candidate kills \(T\): its normalized matrix trace
is faithful on its image, while the displayed squared norm is zero.  In the
target \(\eta=1/2\) quotient, however,
\[
  \operatorname{tr}^{(3)}_{1/2}(T^*T)=\frac1{18}>0,
\]
so \(T\) is nonzero.  In particular, the constructed \(\eta=1/2\) Hecke
representation does not factor through the Temperley--Lieb quotient.

For the remaining case put \(f_i=\Id-e_i\).  The complemented projections
satisfy the same relation with the same \(c\), since
\[
  f_1f_2f_1-f_2f_1f_2
  =-\bigl(e_1e_2e_1-e_2e_1e_2\bigr)=c(f_1-f_2).
\]
Their Markov parameter is \(1-\eta\), because for \(x\in H_n(q)\),
\[
  \operatorname{tr}_\eta(xf_n)
  =\operatorname{tr}_\eta(x)-\operatorname{tr}_\eta(xe_n)
  =(1-\eta)\operatorname{tr}_\eta(x).
\]
At \(\eta=2/3\), therefore, the candidate kills
\[
  T^\perp=f_1f_2f_1-cf_1,
\]
whereas in the target quotient
\[
  \operatorname{tr}^{(3)}_{1/2}\!\left((T^\perp)^*T^\perp\right)
  =\frac1{18}>0.
\]
Neither dimension-three candidate is faithful.  Theorem~\ref{thm:main}
supplies dimension four.
\end{proof}

There is also an immediate infinite family.  For vector spaces \(V,W\), let
\(\Sigma_{23}:(V\otimes W)^{\otimes2}\to V^{\otimes2}\otimes W^{\otimes2}\)
exchange the middle tensor factors.  For
\(R\in\End(V^{\otimes2})\) and \(S\in\End(W^{\otimes2})\), define the
reshuffled tensor product
\[
  R\boxt S=\Sigma_{23}^{-1}(R\otimes S)\Sigma_{23}.
\]
Taking \(S=\Id_{W^{\otimes2}}\) with \(W=\bbC^m\) gives a unitary
Yang--Baxter operator on the base space \(V\otimes W\), of dimension
\(4m\)~\cite[\S3]{Lechner2026}.  Its two eigenvalues and normalized spectral
rank are unchanged.  Therefore
\begin{equation}
  \left[e^{i\pi/3},\frac12,4m\right]\ne\varnothing
  \qquad(m\ge1).
  \label{eq:amplification}
\end{equation}
This does not settle the remaining base dimensions
\(6,10,14,\ldots\).

\section{A (3,2)-generalized normal form}\label{sec:gybe}

The visible identity in the second qubit of every word in~\eqref{eq:H}
does have a precise chain-level interpretation: after a within-site coordinate
swap the chain-level structure becomes explicit.  We use the generalized braid placement of
\cite[Eq.~(3.9) and Remark~3.10(2)]{GHR2013} and the associated tower and
localization conventions of
\cite[Example~4.23 and Definition~4.24]{GHR2013}.  Write the \(i\)-th
four-dimensional site as
\[
  V_i=\bbC^2_{a_i}\otimes\bbC^2_{b_i},
\]
let \(\Sigma_{\mathrm q}(u\otimes v)=v\otimes u\) be the within-site qubit
swap, and put
\[
  U_n=\Sigma_{\mathrm q}^{\otimes n}.
\]
In the factor order obtained after applying \(U_n\),
\[
  (b_1,a_1,b_2,a_2,\ldots),
\]
define
\[
  K=\frac{q-1}{2}\Id_8+\frac{q+1}{2}K_H
\]
and
\begin{equation}
  K_H=
  -\frac{ZZZ+ZJJ+JJZ-JZJ}{\sqrt6}
  -\frac{XXX}{\sqrt3}.
  \label{eq:KH}
\end{equation}
Let \(\rho_R^{(n)}\) denote the standard braid representation generated by
the adjacent copies of \(R\).

\begin{proposition}\label{prop:gybe}
The two-site operator factors after the sitewise swap as
\[
  U_2RU_2^*
  =
  \Id_{b_1}\otimes K_{a_1,b_2,a_2}.
\]
The operator \(K\in\End((\bbC^2)^{\otimes3})\) is unitary and satisfies the
\((3,2)\)-generalized Yang--Baxter equation
\begin{equation}
  (K\otimes\Id_4)(\Id_4\otimes K)(K\otimes\Id_4)
  =
  (\Id_4\otimes K)(K\otimes\Id_4)(\Id_4\otimes K)
  \label{eq:32gybe}
\end{equation}
on \((\bbC^2)^{\otimes5}\), together with far commutativity.  If
\(\rho_K^{(n)}\) denotes the resulting \((3,2)\)-generalized braid
representation, explicitly
\[
  \rho_K^{(n)}(\sigma_i)
  =
  \Id_2^{\otimes2(i-1)}\otimes K
  \otimes\Id_2^{\otimes2(n-i-1)},
\]
then for every \(n\) and every \(x\in\bbC B_n\),
\[
  U_n\rho_R^{(n)}(x)U_n^*
  =
  \Id_{b_1}\otimes\rho_K^{(n)}(x).
\]
Moreover, for the standard inclusion \(\iota_n:\bbC B_n\hookrightarrow
\bbC B_{n+1}\), the placement formula gives
\[
  \rho_K^{(n+1)}(\iota_n(x))
  =\rho_K^{(n)}(x)\otimes\Id_4.
\]
Consequently \(K\) gives a faithful unitary \((3,2)\)-localization of the
tower \(\{H_n(3,6)\}_{n\ge2}\).
\end{proposition}

\begin{proof}
Conjugating the five words in~\eqref{eq:H} by \(U_2\) puts the identity
factor first and leaves exactly~\eqref{eq:KH} on
\((a_1,b_2,a_2)\).  This proves the factorization and unitarity.
On a chain, adjacent copies of \(K\) act on triples shifted by two qubits.
After factoring out the global spectator \(b_1\), the ordinary
Yang--Baxter equation~\eqref{eq:ybe} is precisely
\eqref{eq:32gybe}; nonadjacent triples are disjoint, so far commutativity
holds.  The displayed chain conjugacy follows generator by generator.
The same placement formula gives the displayed compatibility with
\(\iota_n\).
Tensoring by a spectator identity does not change the kernel, and
Proposition~\ref{prop:markov} proves faithfulness on \(H_n(3,6)\).
\end{proof}

Equivalently, any \((3,2)\)-generalized operator \(K\) on \(W^{\otimes3}\)
blocks into an ordinary operator \(\Id_W\otimes K\) on the local space
\(W\otimes W\), in the displayed ordering; the ordinary braid relation is
\(\Id_W\) tensored with~\eqref{eq:32gybe}.  This elementary blocking
observation explains why the present \(8\times8\) active operator yields a
\(16\times16\) ordinary one.

There is an exact tensor-structural distinction from the known operator
\(K_{\mathrm{GHR}}^{\mathrm{gen}}\) in
\cite[Eq.~(5.2), Theorem~5.28, and Remark~5.29(1)]{GHR2013}.  For an
operator \(L\) on three qubits, define
\[
\begin{aligned}
  \mathcal Y_m(L)={}&
  (L\otimes\Id_{2^m})(\Id_{2^m}\otimes L)(L\otimes\Id_{2^m})\\
  &-(\Id_{2^m}\otimes L)(L\otimes\Id_{2^m})(\Id_{2^m}\otimes L).
\end{aligned}
\]
Direct exact calculation gives, with \(\|\cdot\|_F\) the Frobenius norm,
\[
\begin{array}{c@{\qquad}cc}
\toprule
 & \|\mathcal Y_1(\,\cdot\,)\|_F^2
 & \|\mathcal Y_2(\,\cdot\,)\|_F^2\\
\midrule
K_{\mathrm{GHR}}^{\mathrm{gen}} & 0 & 48\\
K                  & 24 & 0\\
\bottomrule
\end{array}
\]
The two nonzero values occur in different ambient matrix dimensions; only
their vanishing or nonvanishing is used here.
Thus the displayed GHR operator is of type \((3,1)\), while the present
operator is of type \((3,2)\).  The distinction concerns tensor placement,
not the two matrices as bare normal operators: they have the same two
eigenvalues with the same multiplicities and hence are unitarily similar
after forgetting tensor structure.  The \((3,2)\) shift is what permits the
spectator blocking \(\Id_2\otimes K\) into an ordinary operator on the local
space \(\bbC^2\otimes\bbC^2\); the GHR operator does not yield this ordinary
blocking in its displayed \((3,1)\) tensor structure.

GHR also records a different \((3,2)\) operator in Example~3.14.  Its two
eigenvalues have projective ratio \(i\) (up to inversion), whereas the ratio
for \(K\) is \(-q=e^{-2\pi i/3}\) (up to inversion), so it is a different
spectral class.  In particular, being of type \((3,2)\) is not by itself the
distinctive feature here; the relevant structure is the exceptional Hecke
spectrum together with a faithful localization of the \(H_n(3,6)\)
Jones--Wenzl tower.

\section{Quaternionic factorization and concurrent comparison}
\label{sec:quaternionic}

The Pauli--Clifford reflection admits an intrinsic quaternionic factorization.
This both explains its relation to Family~III in the concurrent work of
Galindo and Rowell and gives an exact bridge between the two explicit local
formulas.  To distinguish the operator in this section from the
eight-dimensional generalized operator \(K_{\mathrm{GHR}}^{\mathrm{gen}}\)
above, write
\(R_{\mathrm K}:=R\).

\begin{proposition}[Intrinsic Family III form]\label{prop:family-three}
With \(M,E\) as in Lemma~\ref{lem:reflection}, define
\[
  \mathsf A=-i\sqrt2\,M,\qquad \mathsf B=iE,
\]
and
\[
  U_{\mathrm K}=\frac{\mathsf A+\mathsf A\mathsf B}{2},
  \qquad
  V_{\mathrm K}=\frac{\mathsf A-\mathsf A\mathsf B}{2}.
\]
Then
\[
  U_{\mathrm K}^2=V_{\mathrm K}^2=-\Id_{16},\qquad
  U_{\mathrm K}V_{\mathrm K}=-V_{\mathrm K}U_{\mathrm K}=\mathsf B,
\]
and \(U_{\mathrm K},V_{\mathrm K}\) are skew-Hermitian unitaries:
\[
  U_{\mathrm K}^\dagger=-U_{\mathrm K},\qquad
  V_{\mathrm K}^\dagger=-V_{\mathrm K},\qquad
  U_{\mathrm K}^\dagger U_{\mathrm K}
  =V_{\mathrm K}^\dagger V_{\mathrm K}=\Id_{16}.
\]
Moreover,
\begin{equation}
  U_{\mathrm K}+V_{\mathrm K}+U_{\mathrm K}V_{\mathrm K}
  =-i\sqrt3\,H.
  \label{eq:quaternionic-sum}
\end{equation}
If \(\zeta=e^{\pi i/6}\), so that \(q=\zeta^2\), then
\begin{equation}
  R_{\mathrm K}
  =\frac{i\zeta}{2}
   \left(\Id_{16}+U_{\mathrm K}+V_{\mathrm K}
   +U_{\mathrm K}V_{\mathrm K}\right).
  \label{eq:family-three}
\end{equation}
Thus the five-word representative has the abstract quaternionic Family~III
form used in \cite[\S13]{GalindoRowell2026}.
\end{proposition}

\begin{proof}
Lemma~\ref{lem:reflection} gives
\[
  \mathsf A^2=-2\Id,\qquad \mathsf B^2=-\Id,\qquad
  \mathsf A\mathsf B=-\mathsf B\mathsf A.
\]
Write \(U_{\mathrm K}=\frac12\mathsf A(\Id+\mathsf B)\) and
\(V_{\mathrm K}=\frac12\mathsf A(\Id-\mathsf B)\).  Since
\((\Id\pm\mathsf B)\mathsf A=\mathsf A(\Id\mp\mathsf B)\), direct
multiplication gives the two squares and
\(U_{\mathrm K}V_{\mathrm K}=\mathsf B=-V_{\mathrm K}U_{\mathrm K}\).
Since \(M\) and \(E\) are Hermitian, \(\mathsf A\) and \(\mathsf B\) are
skew-Hermitian.  The displayed formulas for \(U_{\mathrm K},V_{\mathrm K}\),
together with \(\mathsf A\mathsf B=-\mathsf B\mathsf A\), show that they too
are skew-Hermitian; their squares then prove unitarity.
Moreover,
\[
  U_{\mathrm K}+V_{\mathrm K}+U_{\mathrm K}V_{\mathrm K}
  =\mathsf A+\mathsf B
  =-i\sqrt2\,M+iE=-i\sqrt3\,H,
\]
where the last equality uses \(H=\sqrt{2/3}\,M-E/\sqrt3\).  Finally
\(i\zeta=q-1\) and \(\sqrt3\,\zeta=q+1\), so
substituting~\eqref{eq:quaternionic-sum} into~\eqref{eq:family-three} yields
Equation~\eqref{eq:R}.
\end{proof}

We now encode the Galindo--Rowell local operator from their literal tensor
placements.  On \(W=\bbC^2\otimes\bbC^2\), set
\[
  P_Z=(I\otimes Z)\otimes(Z\otimes I),\qquad
  P_X=(X\otimes I)\otimes(X\otimes X),
\]
\begin{equation}
  U=iP_Z,\qquad V=iP_X,\qquad
  R_{\mathrm{GR}}=\frac{i\zeta}{2}
  \left(\Id_{16}+U+V+UV\right).
  \label{eq:RGR}
\end{equation}
These are the formulas in \cite[\S13]{GalindoRowell2026}, not the older
\(8\times8\) operator \(K_{\mathrm{GHR}}^{\mathrm{gen}}\) from GHR
Equation~(5.2).
Let \(\Sigma(x\otimes y)=y\otimes x\) swap the two four-dimensional sites.

\begin{proposition}[Exact comparison with Galindo--Rowell]
\label{prop:concurrent-comparison}
Let \(\varepsilon_0,\varepsilon_1\) be the standard basis of \(\bbC^2\).  In
the ordered basis
\((\varepsilon_0\otimes\varepsilon_0,\varepsilon_0\otimes\varepsilon_1,
\varepsilon_1\otimes\varepsilon_0,\varepsilon_1\otimes\varepsilon_1)\) of \(W\),
the matrix
\[
S=\frac14\begin{pmatrix}
2+\sqrt2 & -i\sqrt2 & i\sqrt2 & 2-\sqrt2\\
\sqrt2 & i(2+\sqrt2) & i(2-\sqrt2) & -\sqrt2\\
-(2-\sqrt2) & -i\sqrt2 & i\sqrt2 & -(2+\sqrt2)\\
-\sqrt2 & i(2-\sqrt2) & i(2+\sqrt2) & \sqrt2
\end{pmatrix}
\]
is unitary and satisfies
\begin{equation}
  R_{\mathrm K}
  =(S^\dagger\otimes S^\dagger)\,
    \Sigma R_{\mathrm{GR}}\Sigma\,(S\otimes S).
  \label{eq:concurrent-comparison}
\end{equation}
Hence \(R_{\mathrm K}\) is a local-unitary normal form of the opposite
operator \(R_{\mathrm{GR},21}:=\Sigma R_{\mathrm{GR}}\Sigma\).
\end{proposition}

\begin{proof}
Exact row inner products give \(SS^\dagger=\Id_4\), hence
\(S^\dagger S=\Id_4\).  For
\[
  \mathcal C(T)=(S^\dagger\otimes S^\dagger)\Sigma T\Sigma(S\otimes S),
\]
substitution of the displayed matrices gives
\[
  \mathcal C(U)=U_{\mathrm K},\qquad
  \mathcal C(V)=V_{\mathrm K},\qquad
  \mathcal C(UV)=U_{\mathrm K}V_{\mathrm K}.
\]
Applying \(\mathcal C\) to~\eqref{eq:RGR} and using
Proposition~\ref{prop:family-three} proves~\eqref{eq:concurrent-comparison}.
All entries in these finite products lie in
\(\bbQ(\sqrt2,\sqrt3,i)\); the independent exact verifier in the source
package checks the three generator identities as well as the resulting
operator identity.
\end{proof}

Equation~\eqref{eq:concurrent-comparison} includes a site reversal.  The
transformation \(R\mapsto\Sigma R\Sigma\) preserves the Yang--Baxter equation,
so the opposite is again a Yang--Baxter operator.  Galindo--Rowell's stated
local-equivalence relation permits a unit scalar and a common local basis
change~\cite[\S2]{GalindoRowell2026}, but does not include this site reversal.
Proposition~\ref{prop:concurrent-comparison} therefore proves local-unitary
equivalence to the opposite operator \(R_{\mathrm{GR},21}\).  It does not
determine whether a direct local-unitary equivalence to \(R_{\mathrm{GR}}\)
also exists.  The two formulas represent the same underlying local solution up
to opposite and local basis change.  This comparison is separate from the
preceding residual table for \(K_{\mathrm{GHR}}^{\mathrm{gen}}\), which
distinguishes the displayed \((3,1)\) and \((3,2)\) generalized tensor
placements.

\section{Finite braid images and the enhanced link invariant}
\label{sec:global}

We now pass from the exact two-site comparison to the full braid
representations and to the link invariant in the normalization of the
five-word matrix.  Put
\[
  \kappa=q-1=q^2=e^{2\pi i/3}.
\]
Propositions~\ref{prop:enhancement}--\ref{prop:twists} are direct calculations
from the displayed five-word operator and its scalar partial traces.
Theorem~\ref{thm:all-strand} promotes the local opposite comparison to all
strands.  Corollaries~\ref{cor:clifford}--\ref{cor:algorithm} transfer the known
Family~III consequences, while Corollary~\ref{cor:branched-cover} records the
classical Lickorish--Millett identification.

\begin{proposition}[Scalar enhancement]\label{prop:enhancement}
The quadruple
\[
  (R,\mu,\alpha,\beta)=(R,\Id_4,\kappa,2)
\]
is an enhanced unitary Yang--Baxter operator in Turaev's convention.  Its
oriented-link invariant is
\begin{equation}
  \mathcal J_R(\widehat\xi)
  =\kappa^{-\operatorname{wr}(\xi)}2^{-n}
   \Tr(\rho_n(\xi)),
  \qquad \xi\in B_n,
  \label{eq:enhanced-invariant}
\end{equation}
where \(\Tr\) is the ordinary unnormalized trace on the \(4^n\)-dimensional
tensor space and \(\operatorname{wr}\) is the braid exponent sum.  In
particular, \(\mathcal J_R(\bigcirc)=2\).  Since \(\mu=\Id_4\), the factor
\(\mu^{\otimes n}\) in Turaev's general trace formula is \(\Id_{4^n}\) and is
suppressed in~\eqref{eq:enhanced-invariant}.
\end{proposition}

\begin{proof}
Equations~\eqref{eq:R} and~\eqref{eq:H-identities}, together with the partial
traces computed in the proof of Theorem~\ref{thm:main}, give
\[
  \Tr_1(R)=\Tr_2(R)=2\kappa\Id_4.
\]
The spectral decomposition also gives
\[
  R^{-1}=\frac{q^{-1}-1}{2}\Id_{16}
         +\frac{q^{-1}+1}{2}H.
\]
Since \(q^{-1}-1=\kappa^{-1}\),
\[
  \Tr_1(R^{-1})=\Tr_2(R^{-1})=2\kappa^{-1}\Id_4.
\]
Turaev's enhancement conditions are
\[
  R(\mu\otimes\mu)=(\mu\otimes\mu)R,
\]
\[
  \Tr_2(R(\mu\otimes\mu))=\alpha\beta\mu,
  \qquad
  \Tr_2(R^{-1}(\mu\otimes\mu))=\alpha^{-1}\beta\mu
\]
\cite[\S\S2.2--2.3]{Turaev1988}.  They hold with the displayed scalars; the
commutation is immediate.  Turaev's trace formula and Markov theorem then give
\eqref{eq:enhanced-invariant}
\cite[\S\S3.1--3.2, especially Theorem~3.1.2]{Turaev1988}.  For the closure of the identity in
\(B_1\), this is \(2^{-1}\Tr(\Id_4)=2\).
\end{proof}

\begin{proposition}[HOMFLYPT specialization]\label{prop:homflypt}
Let \(P_{\mathrm H}\) be the normalized HOMFLYPT polynomial in the convention
\[
  aP_{\mathrm H}(L_+)-a^{-1}P_{\mathrm H}(L_-)
  =zP_{\mathrm H}(L_0),
  \qquad P_{\mathrm H}(\bigcirc)=1.
\]
Then
\begin{equation}
  \mathcal J_R(L)=2P_{\mathrm H}(L;i,i).
  \label{eq:homflypt-specialization}
\end{equation}
Consequently \(\mathcal J_R\) is invariant under mirror image, and the
\(c\)-component unlink has value \(2^c\).
\end{proposition}

\begin{proof}
Multiplying the Hecke relation by \(R^{-1}\) yields
\begin{equation}
  R-qR^{-1}=(q-1)\Id_{16}=\kappa\Id_{16}.
  \label{eq:matrix-skein}
\end{equation}
Apply ordinary trace to braid closures that differ only by \(R\), \(R^{-1}\),
or the identity at one crossing.  Accounting for their writhes in
\eqref{eq:enhanced-invariant} gives
\[
  \mathcal J_R(L_+)-q\kappa^{-2}\mathcal J_R(L_-)
  =\mathcal J_R(L_0).
\]
Because \(\kappa=q^2\) and \(q^3=-1\),
\[
  q\kappa^{-2}=q^{-3}=-1,
\]
and hence
\begin{equation}
  \mathcal J_R(L_+)+\mathcal J_R(L_-)=\mathcal J_R(L_0).
  \label{eq:plus-skein}
\end{equation}
At \((a,z)=(i,i)\), the displayed HOMFLYPT relation is exactly the same
relation.  Its normalized skein characterization and
\(\mathcal J_R(\bigcirc)=2\) prove~\eqref{eq:homflypt-specialization}.  The
relation~\eqref{eq:plus-skein} is symmetric under interchange of positive and
negative crossings, giving mirror invariance.  The unlink value follows either
from the standard split-union factor or directly from
\eqref{eq:enhanced-invariant}.
\end{proof}

\begin{proposition}[Local order and twist periodicity]\label{prop:twists}
One has
\[
  R^3=-\Id_{16},\qquad R^6=\Id_{16}.
\]
If \(L_m=\widehat{x\sigma_i^m y}\), with the other braid factors fixed, then
\[
  \mathcal J_R(L_{m+3})=-\mathcal J_R(L_m),\qquad
  \mathcal J_R(L_{m+6})=\mathcal J_R(L_m).
\]
\end{proposition}

\begin{proof}
The eigenvalues of \(R\) are \(-1,q\), and both have cube \(-1\).  This proves
the two operator identities.  Replacing \(m\) by \(m+3\) multiplies the
represented braid by \(-\Id\) and increases writhe by three.  Since
\(\kappa^3=1\), the writhe factor is unchanged, proving the first formula; the
second follows by iteration.  This is periodicity within a local twist family,
not an identification of the full braid image group.
\end{proof}

\begin{theorem}[All-strand equivalence]\label{thm:all-strand}
Let \(\rho_n^{\mathrm K}\) and \(\rho_n^{\mathrm{GR}}\) denote the braid
representations generated by \(R_{\mathrm K}\) and \(R_{\mathrm{GR}}\),
respectively, and put \(S_n=S^{\otimes n}\).  Let
\(\mathsf{Rev}_n\) reverse the order of the \(n\) four-dimensional tensor
sites, and let \(\theta_n\in\operatorname{Aut}(B_n)\) be defined by
\(\theta_n(\sigma_i)=\sigma_{n-i}\).  Then
\begin{equation}
  \mathsf{Rev}_nS_n\rho_n^{\mathrm K}(\xi)
  S_n^\dagger\mathsf{Rev}_n^\dagger
  =\rho_n^{\mathrm{GR}}(\theta_n(\xi))
  \qquad(\xi\in B_n).
  \label{eq:reflected-global}
\end{equation}
Moreover, for the Garside half twist
\[
  \Delta_n=(\sigma_1\cdots\sigma_{n-1})
  (\sigma_1\cdots\sigma_{n-2})\cdots\sigma_1,
\]
put
\[
  D_n=\rho_n^{\mathrm{GR}}(\Delta_n),\qquad
  \mathcal U_n=D_n^\dagger\mathsf{Rev}_nS_n.
\]
Then
\begin{equation}
  \mathcal U_n\rho_n^{\mathrm K}(\xi)\mathcal U_n^\dagger
  =\rho_n^{\mathrm{GR}}(\xi)
  \qquad(\xi\in B_n).
  \label{eq:same-word-global}
\end{equation}
Thus the image groups are unitarily conjugate and ordinary traces agree on the
same braid word.
\end{theorem}

\begin{proof}
Equation~\eqref{eq:concurrent-comparison} is equivalent to
\[
  (S\otimes S)R_{\mathrm K}(S^\dagger\otimes S^\dagger)
  =\Sigma R_{\mathrm{GR}}\Sigma.
\]
Conjugating a two-site block by the global site reversal reverses its two
internal sites and moves positions \((i,i+1)\) to \((n-i,n-i+1)\).  The two
reversals cancel, so generator by generator
\[
  \mathsf{Rev}_nS_n\rho_n^{\mathrm K}(\sigma_i)
  S_n^\dagger\mathsf{Rev}_n^\dagger
  =\rho_n^{\mathrm{GR}}(\sigma_{n-i}).
\]
This proves~\eqref{eq:reflected-global}.  The standard half-twist identity
\[
  \Delta_n\sigma_i\Delta_n^{-1}=\sigma_{n-i}
\]
says \(\theta_n=\operatorname{Ad}_{\Delta_n}\).  Hence the right side of
\eqref{eq:reflected-global} is
\(D_n\rho_n^{\mathrm{GR}}(\xi)D_n^\dagger\); conjugating by \(D_n^\dagger\)
proves~\eqref{eq:same-word-global}.  The local comparison includes an opposite,
but the Garside conjugation removes the reflected generator indices at the
level of the complete \(n\)-strand representation.  It does not establish a
direct two-site local equivalence without the opposite.
\end{proof}

\begin{corollary}[Finite images and a Clifford frame]\label{cor:clifford}
For every \(n\ge2\), the group \(\rho_n^{\mathrm K}(B_n)\) is finite.  It is
Clifford with respect to a fixed conjugated Pauli frame.  In particular, the
images are not dense in the ambient unitary groups and are not computationally
universal by braiding alone in the usual density-based sense.
\end{corollary}

\begin{proof}
In the Galindo--Rowell realization, \(U=iP_Z\) and \(V=iP_X\) for the Hermitian
Pauli words in~\eqref{eq:RGR}.  Since \(i\zeta=\kappa\),
\[
  R_{\mathrm{GR}}
  =\kappa\left(\frac{\Id+U}{\sqrt2}\right)
             \left(\frac{\Id+V}{\sqrt2}\right)
  =\kappa e^{i\pi P_Z/4}e^{i\pi P_X/4}.
\]
Each factor is a Pauli quarter-turn and hence normalizes the Pauli group; this
is the Family~III Clifford structure of
\cite[Proposition~8.1]{GalindoRowell2026}.  If \(\mathcal P_{2n}\) is the
standard Pauli group on the \(2n\) qubits, then the original generators
normalize the uniformly conjugated frame
\[
  (\mathsf{Rev}_nS_n)^\dagger\mathcal P_{2n}(\mathsf{Rev}_nS_n).
\]
The need for the conjugated frame is concrete.  Using the standard Hermitian
Pauli \(Y=iJ\) and writing \(XIII=X\otimes I\otimes I\otimes I\), exact
conjugation gives
\[
\begin{aligned}
  R_{\mathrm K}(XIII)R_{\mathrm K}^\dagger
  =\frac1{2\sqrt2}\bigl(&YIYY-YIYZ+YIZY-YIZZ\\
                         &{}-ZIYY+ZIYZ-ZIZY+ZIZZ\bigr).
\end{aligned}
\]
The eight Pauli words are distinct, so \(R_{\mathrm K}\) does not normalize the
standard computational Pauli group; its Clifford statement is specifically
relative to the displayed conjugated frame.  No assertion that \(S\) itself is
a standard computational Clifford operator is needed.  Finiteness of the
quaternionic braid image is
\cite[Theorem~3.1]{Rowell2011}; after the finite-order scalar normalization, it
also follows in the broader Family~III framework from
\cite[Theorem~8.4]{GalindoRowell2026}.  Theorem~\ref{thm:all-strand} transfers it
to \(\rho_n^{\mathrm K}\).
\end{proof}

\begin{corollary}[Quaternionic tower]\label{cor:quaternionic-tower}
Under the all-strand equivalence, the localized quotient is the
braid-generated subalgebra of the explicit quaternionic twisted-group-algebra
tower, where \(s_i\) denotes Rowell's quaternionic braid generator:
\[
  H_n(3,6)\cong
  \langle s_1,\ldots,s_{n-1}\rangle_{\mathrm{alg}}
  \subset \mathcal A_n(Q_8).
\]
The full tower algebra has the normal basis
\[
  u_1^{\epsilon_1}v_1^{\delta_1}\cdots
  u_{n-1}^{\epsilon_{n-1}}v_{n-1}^{\delta_{n-1}},
  \qquad \epsilon_i,\delta_i\in\{0,1\},
\]
and
\[
  \dim_{\bbC}\mathcal A_n(Q_8)=4^{n-1}.
\]
\end{corollary}

\begin{proof}
The basis and dimension are the literal construction in
\cite[\S13.1]{GalindoRowell2026}.  The identification of the braid-generated
subalgebra with \(H_n(3,6)\) is
\cite[Lemma~3.2, p.~176]{Rowell2011}.  This does not identify
\(H_n(3,6)\) with the full \(4^{n-1}\)-dimensional tower algebra.
\end{proof}

\begin{corollary}[Exact classical evaluation]\label{cor:algorithm}
Given \(\xi\in B_n\), the value \(\mathcal J_R(\widehat\xi)\) is exactly
computable by a deterministic classical algorithm in time polynomial in
\(n+|\xi|\).
\end{corollary}

\begin{proof}
For Family~III with the fixed metric group \(A=\mathbb Z/2\mathbb Z\),
\cite[Theorem~7.10]{GalindoRowell2026} gives an exact deterministic algorithm,
polynomial in strand number plus braid-word length, for the coefficient-trace
invariant.  To match normalizations, put
\[
  r_0=\kappa^{-1}R_{\mathrm{GR}}
  =\frac12(\Id+U+V+UV).
\]
If \(\varepsilon_n\) is their identity-coefficient functional and \(\Phi_n\)
the Pauli realization, then \cite[Theorem~7.3]{GalindoRowell2026} gives
\[
  \Tr\!\circ\!\Phi_n=4^n\varepsilon_n.
\]
Because \(\rho_{R_{\mathrm{GR}}}(\xi)
=\kappa^{\operatorname{wr}(\xi)}\Phi_n(\rho_{r_0}(\xi))\), its enhanced trace
is \(2^n\varepsilon_n(\rho_{r_0}(\xi))\), exactly the invariant covered by
that theorem.  The same-word equivalence~\eqref{eq:same-word-global} transfers
the algorithm to \(R_{\mathrm K}\).  This complexity statement comes from the
cited algorithm, not from the exponentially growing dimension in
Corollary~\ref{cor:quaternionic-tower}.
\end{proof}

\begin{corollary}[Triple cyclic branched cover]\label{cor:branched-cover}
Let \(L\) be an oriented link with \(c(L)\) components.  Let
\(\Sigma_3(L)\) be the connected triple cyclic cover of \(S^3\) branched over
\(L\), obtained by completing the cover of \(S^3\setminus L\) associated with
the homomorphism that sends every positively oriented meridian to
\(1\in\mathbb Z/3\mathbb Z\).  Put
\[
  d_2(L)=\dim_{\mathbb F_2}H_1(\Sigma_3(L);\mathbb F_2).
\]
Then \(d_2(L)\) is even and
\begin{equation}
  \mathcal J_R(L)=2(-1)^{c(L)-1}(-2)^{d_2(L)/2}.
  \label{eq:branched-cover}
\end{equation}
\end{corollary}

\begin{proof}
Lickorish--Millett normalize their polynomial \(P_L(\ell,m)\) by
\[
  P_{\bigcirc}=1,\qquad
  \ell P_{L_+}+\ell^{-1}P_{L_-}+mP_{L_0}=0.
\]
Their Theorem~2, together with the parity-preserving dimension calculation in
its proof, gives
\[
  d_2(L)\in2\mathbb Z,\qquad P_L(1,1)=(-2)^{d_2(L)/2}
\]
\cite[Theorem~2 and pp.~353--355]{LickorishMillett1986}.  In an oriented skein
triple, \(c(L_+)=c(L_-)\), while \(c(L_0)\) differs by one.  Therefore
\[
  Q(L)=(-1)^{c(L)-1}P_L(1,1)
\]
satisfies \(Q(L_+)+Q(L_-)=Q(L_0)\) and \(Q(\bigcirc)=1\).  Comparison with the
normalization in Proposition~\ref{prop:homflypt} gives
\[
  P_{\mathrm H}(L;i,i)=(-1)^{c(L)-1}(-2)^{d_2(L)/2}.
\]
Equation~\eqref{eq:homflypt-specialization} proves~\eqref{eq:branched-cover}.
This is the classical Lickorish--Millett evaluation; the point here is its exact
identification with the invariant of the displayed five-word matrix.
\end{proof}

For orientation and crossing conventions, the exact matrix trace gives the
following small checks:
\[
\begin{array}{c@{\qquad}c}
\toprule
\text{closure} & \mathcal J_R\\
\midrule
1\in B_1\ (\text{unknot}) & 2\\
1\in B_2\ (\text{two-component unlink}) & 4\\
\sigma_1^{\pm2}\in B_2\ (\text{Hopf links}) & -2\\
\sigma_1^{\pm3}\in B_2\ (\text{trefoils}) & -4\\
\bigl(\sigma_1\sigma_2^{-1}\bigr)^2\in B_3\ (\text{figure-eight knot}) & -4\\
\bigl(\sigma_1\sigma_2^{-1}\bigr)^3\in B_3\ (\text{Borromean rings}) & 2\\
\bottomrule
\end{array}
\]
These values agree with~\eqref{eq:branched-cover}; the positive and negative
closures also illustrate mirror insensitivity.  For the \(B_3\) rows,
\(\sigma_1\) and \(\sigma_2\) act by the adjacent placements
\(R\otimes\Id_4\) and \(\Id_4\otimes R\), respectively.

\section{Verification and provenance}\label{sec:verification}

The source package supplies five supported exact-verification routes:
\begin{enumerate}[leftmargin=2.2em,itemsep=0.25em]
\item a hardened SymPy matrix verifier based on the original discovery-era
checker;
\item an independently written standard-library matrix verifier over
\(\bbQ(\sqrt2,\sqrt3,i)\);
\item a matrix-free Pauli-word verifier for the complete \(18\)-word
certificate in Proposition~\ref{prop:circle};
\item an independently encoded verifier for the intrinsic quaternionic form,
the literal Galindo--Rowell operator, the site swap, and \(S\); and
\item a braid-and-link verifier that rechecks the intrinsic factorization,
\(S\)-unitarity, and two-site comparison, then checks the scalar enhancement,
matrix skein identity, local order, two- and three-strand link values, the
standard-frame non-Clifford witness, Pauli quarter-turns, and the global
reversal and Garside conjugacies at \(n=3,4\).
\end{enumerate}
Together these routes check the displayed matrix identities, partial traces,
both dimension-three obstructions, the six-dimensional three-strand Hecke
image, the generalized operator and far commutativity, the exact residual
comparison with \(K_{\mathrm{GHR}}^{\mathrm{gen}}\), and the finite identities
in Sections~\ref{sec:quaternionic}--\ref{sec:global}.  All comparisons are
exact.  Negative tests deliberately alter coefficients, tensor placements,
the local unitary, the quarter-turn order, \(\zeta\), \(\kappa\), the writhe
factor, the skein sign, a reversed generator index, and the Garside word.  The
README records the commands and scope.  These finite checks do not replace the
printed tower-faithfulness and all-strand conjugacy proofs, Lechner's
classification input, or the literature-based interpretation of the
concurrent and topological results.

The five-word support was originally found by an AI-assisted structured
numerical search in a real Pauli--Clifford basis.  The discovery code and
random seeds were not retained.  Accordingly, this paper makes no claim that
the search was reproducible, exhaustive, or established uniqueness.  The
generative-AI systems identified in the declaration below were also used for
mathematical exploration and derivation, adversarial proof analysis, verifier
development, literature support, and manuscript preparation.  Every theorem
claim is supported by the printed derivation, and the finite identities are
checked by the exact programs.  No conclusion depends on accepting model
output as evidence or reproducing the original numerical search.

\section{Related work, novelty, and limitations}\label{sec:priority}

Rowell and Wang introduced the localization framework and conjecture
\cite{RowellWang2012}.  Galindo, Hong, and Rowell then studied the
\(\cC(\mathfrak{sl}_3,6)\) sequence, excluded ordinary local dimension two,
and supplied quasi- and \((3,1)\)-generalized localizations, but no ordinary
dimension-four operator~\cite[\S5.6]{GHR2013}.  Rowell's quaternionic model
and the unitary-braided-vector-space literature are foundational prior art
\cite{Rowell2011,GalindoRowell2014}.  Lechner later identified
\([e^{i\pi/3},1/2,2m]\) as the remaining existence family up to
braid-character equivalence and stated that its dimension-two member is empty
\cite[p.~16, discussion following Proposition~3.6]{Lechner2026}.

Galindo and Rowell independently realized the same exceptional local class and
proved faithful strict localization of \(H_n(3,6)\) in local dimension four,
the smallest dimension in Lechner's exceptional family, using a quaternionic
twisted-group-algebra tower
\cite[\S13]{GalindoRowell2026}.  Their construction is substantially broader
in its classification setting.  The present proof architecture is different:
it gives a five-word real Pauli--Clifford representative, reduces the braid
relation to two anticommuting involutions, prints the complete \(18\)-word
cubic certificate, derives the Markov trace from scalar partial traces, proves
the dimension-three obstruction directly, and exhibits the active
\((3,2)\)-generalized normal form.  Proposition~\ref{prop:concurrent-comparison}
supplies an exact bridge to their Family~III formula, and
Theorem~\ref{thm:all-strand} extends it to every braid word.

The public and reported private chronology is recorded in the Note added in
the Introduction; the papers are treated throughout as independent and
concurrent.  The finite-image, Clifford, twisted-tower, and classical
link-evaluation phenomena are known for the quaternionic representation.  The
contribution here is their exact transfer to, and transparent derivation from,
the sparse five-word representative, including its scalar enhancement and
matrix-trace normalization.  No new HOMFLYPT evaluation, branched-cover
theorem, or all-\(n\) finite-image group classification is claimed.  The result
does not classify the exceptional family in every even dimension or prove
uniqueness.  The comparison theorem concerns the \(16\times16\) ordinary
Family~III operator \(R_{\mathrm{GR}}\), not the older \(8\times8\) generalized
operator \(K_{\mathrm{GHR}}^{\mathrm{gen}}\).

\section*{Data, code, and source availability}

The source, exact verifiers, frozen outputs, checksums, and audit materials for
this edition accompany the manuscript in a deterministic version 1.2.0
package identified by the version-specific Zenodo DOI
\href{https://doi.org/10.5281/zenodo.22013710}
{doi:10.5281/zenodo.22013710}~\cite{Kriebel2026Zenodo}.  The preceding version
1.1.3 package remains permanently archived at
\href{https://doi.org/10.5281/zenodo.21971507}
{doi:10.5281/zenodo.21971507}; it is not overwritten by this revision.  The
maintained \href{https://aleckriebel.github.io/Math/papers/exceptional-ybe-d4/}
{project page} provides the public manuscript and verification links.  No
empirical dataset was generated or analyzed.

\section*{CRediT authorship contribution statement}

Alec Kriebel: Conceptualization, formal analysis, investigation, methodology,
software, validation, writing--original draft, and writing--review and editing.

\section*{Funding}

This research received no external funding.

\section*{Declaration of competing interest}

The author declares no competing interests.

\section*{Declaration of generative AI and AI-assisted technologies in the
manuscript preparation process}

During preparation of this work, the author used OpenAI GPT-5.6 Sol in Pro
mode through ChatGPT at chatgpt.com and OpenAI GPT-5.6 Sol in Ultra mode
through Codex in the ChatGPT desktop application for research support, review,
and composition, and Anthropic Claude through claude.ai for review.  The
author reviewed and edited the content as needed and takes full responsibility
for the content of the publication.

\begin{thebibliography}{10}

\bibitem{GHR2013}
C.~Galindo, S.-M. Hong, and E.~C. Rowell,
\newblock Generalized and quasi-localizations of braid group representations,
\newblock \emph{Int. Math. Res. Not.} \textbf{2013} (2013), no.~3, 693--731.
\newblock \href{https://doi.org/10.1093/imrn/rnr269}{doi:10.1093/imrn/rnr269}.

\bibitem{GalindoRowell2014}
C.~Galindo and E.~C. Rowell,
\newblock Braid representations from unitary braided vector spaces,
\newblock \emph{J. Math. Phys.} \textbf{55} (2014), 061702.
\newblock \href{https://doi.org/10.1063/1.4880196}{doi:10.1063/1.4880196}.

\bibitem{GalindoRowell2026}
C.~Galindo and E.~C. Rowell,
\newblock Unitary Yang--Baxter operators: towards a classification,
\newblock arXiv:2608.16865v1 (17 August 2026).
\newblock \href{https://doi.org/10.48550/arXiv.2608.16865}
{doi:10.48550/arXiv.2608.16865};
\url{https://arxiv.org/abs/2608.16865v1}.

\bibitem{Kriebel2026Zenodo}
A.~Kriebel,
\newblock An exceptional four-dimensional unitary Hecke Yang--Baxter operator:
a five-word Pauli--Clifford normal form,
\newblock Zenodo, version 1.2.0 (2026).
\newblock \href{https://doi.org/10.5281/zenodo.22013710}
{doi:10.5281/zenodo.22013710}.

\bibitem{LickorishMillett1986}
W.~B.~R. Lickorish and K.~C. Millett,
\newblock Some evaluations of link polynomials,
\newblock \emph{Comment. Math. Helv.} \textbf{61} (1986), no.~1, 349--359.
\newblock \href{https://doi.org/10.1007/BF02621920}
{doi:10.1007/BF02621920}.

\bibitem{Lechner2026}
G.~Lechner,
\newblock The classification problem for unitary \(R\)-matrices with two
eigenvalues,
\newblock arXiv:2603.20158v1 (20 March 2026).
\newblock \href{https://doi.org/10.48550/arXiv.2603.20158}
{doi:10.48550/arXiv.2603.20158};
\url{https://arxiv.org/abs/2603.20158v1}.

\bibitem{Rowell2011}
E.~C. Rowell,
\newblock A quaternionic braid representation (after Goldschmidt and Jones),
\newblock \emph{Quantum Topol.} \textbf{2} (2011), 173--182.
\newblock \href{https://doi.org/10.4171/QT/18}{doi:10.4171/QT/18}.

\bibitem{RowellWang2012}
E.~C. Rowell and Z.~Wang,
\newblock Localization of unitary braid group representations,
\newblock \emph{Comm. Math. Phys.} \textbf{311} (2012), 595--615.
\newblock \href{https://doi.org/10.1007/s00220-011-1386-7}
{doi:10.1007/s00220-011-1386-7}.

\bibitem{Turaev1988}
V.~G. Turaev,
\newblock The Yang--Baxter equation and invariants of links,
\newblock \emph{Invent. Math.} \textbf{92} (1988), no.~3, 527--553.
\newblock \href{https://doi.org/10.1007/BF01393746}
{doi:10.1007/BF01393746}.

\bibitem{Wenzl1988}
H.~Wenzl,
\newblock Hecke algebras of type \(A_n\) and subfactors,
\newblock \emph{Invent. Math.} \textbf{92} (1988), 349--383.
\newblock \href{https://doi.org/10.1007/BF01404457}
{doi:10.1007/BF01404457}.

\end{thebibliography}

\end{document}
